\section{Current Evaluation Protocol and Evidence}
\label{app:current-evaluation}
\paperparagraph{Coverage and acceptance.}
We evaluate all 123 modeled cases with one optimized compiler under
algebraic semantics and a 30-second per-command limit. Acceptance requires
semantic checking, successful compilation, and portable-program validation
under the original guards and supplied fields. Coverage uses the compiler's
\code{check}, \code{compile}, and \code{validate} commands in sequence.
The 86 unmodeled catalog cases are excluded from
semantic evaluation and counted as neither successes nor failures.
Exact sources, command invocations, per-case outcomes, and compiler provenance are retained in
\nolinkurl{evaluation/unified-construction/}.

\begin{table}[!htbp]
\caption{Selected operation families in the coverage run.
\textbf{Memberships overlap; do not sum rows.} Entries count all modeled cases (axis-inference cases).}
\label{tab:current-operation-families}
\small
\begin{tabularx}{\linewidth}{@{}Yrr@{}}\toprule
Family & Cases (axis) & Accepted (axis)\\\midrule
Dot & 12 (12) & 11 (11)\\
Transpose & 10 (8) & 10 (8)\\
Reduction & 9 (9) & 9 (9)\\
Addition & 13 (3) & 12 (2)\\
Broadcast & 19 (19) & 15 (15)\\
Reverse & 10 (8) & 7 (6)\\
Slice & 10 (5) & 2 (0)\\
Convolution & 1 (1) & 0 (0)\\\bottomrule
\end{tabularx}
\end{table}
The public compile/load/rewrite suite passes 685 independent dense-value
and ordered-shape checks and 142 rejection checks across eight test targets.
It exercises degenerate extents, varying ranks, axis ordering, and malformed
programs. Shared-IR regressions additionally check native/symbolic agreement,
residual-program retention, and paired contraction reindexing. Detailed
inventories accompany the artifacts; these development tests and corpus
reruns do not constitute a new held-out evaluation.

\paperparagraph{Specification accounting.}
The source-footprint comparison includes all 31 rule identities
whose axis inference succeeds in the current coverage run. We recount the
saved original explicit and compact inferred sources from the earlier
32-case size protocol, checking their source hashes and recorded token
counts. No new compiler experiment is performed. Five original explicit
forms and one compact form did not validate in that historical size run;
all 31 identities remain included in this static source comparison.
The current coverage result concerns the axis-inference cases and does not
establish acceptance of every counted source variant.

The original explicit forms contain 3,014 tokens and the compact inferred
forms contain 2,559, a difference of 455 tokens (15.1\%). Saved inferred
forms include unit-axis normalization and unused-label/empty-block cleanup.
Two normalized cases add 48 legality tokens relative to the originals, so
the overall difference is not attributed solely to inference. Exact counted
sources, per-rule responsibility counts, historical outcomes, and source
hashes are retained in
\nolinkurl{evaluation/current/specification_size/results/all_31_source_footprint.json}.

Table~\ref{tab:xla-specification-size}'s separate XLA comparison covers
25 audited cases, with 9,188 XLA, 2,224 explicit TEPL, and 1,829 inferred
TEPL tokens. These historical measurements are retained in
\nolinkurl{evaluation/current/specification_size/results/responsibility.json}.
Some selected XLA branches cover additional cases and policies, so this
comparison does not establish equal implementation scope.

Lexical tokens exclude comments and whitespace; imports and shared
operation declarations are excluded. The five responsibilities partition
each included token exactly once. Shared dialect and compiler costs remain
real reusable costs outside these per-rule counts. Supplement S2 gives the
worked transpose example.

\FloatBarrier
\sidefigure[0.6]{figures/scalability.pdf}
{Certified compilation scalability. (a) Compilation time per phase grows approximately linearly with dot-chain depth. (b) Compilation time remains stable across the tested tensor ranks.}
{fig:scalability}
{Two horizontally aligned panels show certified dot-chain compilation with stacked mean phase costs and overlaid median total latency, and transpose-rank compilation. Whiskers show minimum and maximum total latency. All fifteen automatic configurations succeed in every measured round.}

\paperparagraph{Compilation scalability.}
We evaluate the scalability of the core algorithm as rewrite depth and
tensor rank increase. Figure~\ref{fig:scalability} shows compilation time
for dot-chain reassociation and transpose composition, including the
dot-chain cost by compilation phase.

The experiment uses 1--10 dots per side at input rank three with fixed
source axis lists, and three source transposes composed into one replacement
at ranks 2, 4, 8, 16, and 32 with symbolic permutations. All RHS axis fields
are inferred. We run the optimized compiler with guard certification under
algebraic semantics, using one warmup and ten measured rounds per
configuration. Timing covers compilation through executable validation;
parsing and import loading are excluded. Supplement~\cite{tepl-supplement},
Section~S3, gives the complete protocol, phase boundaries, and timing table.

All 15 configurations compile and validate in every measured round.
Dot-chain compilation time increases approximately linearly with chain
depth, from a median of 1.49 ms for one dot per side to 6.06 ms for ten
(Figure~\ref{fig:scalability}(a)). Tensor rank has little effect on
transpose compilation time over the tested range: medians remain between
1.16 and 1.32 ms, with no systematic increase from rank 2 to rank 32
(Figure~\ref{fig:scalability}(b)).
